\chapter*{Prefaci} \addcontentsline{toc}{chapter}{Prefaci}

   Voldria dir en primer lloc que no he intentat escriure un tractat complet
   d'electrotècnia, d'electrònica o de sistemes elèctrics de potència, sinó que més aviat
   he volgut
   fer una recopilació de qüestions teòriques i pràctiques relacionades amb els camps mencionats
   anteriorment.

   Les fonts d'informació utilitzades en la realització d'aquest llibre són molt diverses,
   i inclouen llibres sobre les diferents matèries, articles de revistes o d'Internet,
   apunts de classe d'assignatures impartides a l'ETSEIB,\index{ETSEIB} i d'altres.

   Pel que fa al llibre en si mateix, ha estat escrit utilitzant el sistema de composició de
   textos \LaTeX,\index{LaTex@\LaTeX} el qual
   permet integrar molt fàcilment text, fórmules i gràfics, obtenint un resultat de
   gran qualitat. S'han utilitzat diversos paquets d'ampliació, com ara
   l'\AmS-\LaTeX,\index{AmSLaTex@\AmS-\LaTeX}
   per tal de millorar la presentació de certes parts del
   llibre, com per exemple les taules, les capçaleres o les fórmules matemàtiques. S'ha utilitzat la distribució MiK\TeX, que ofereix una implementació lliure de \LaTeX , accessible a l'adreça: \href{http://www.miktex.org/}{www.miktex.org}.

   Aquest llibre està pensat per ser utilitzat tant de forma directa en la pantalla d'un
   ordinador com en paper després d'haver-lo imprès; per tal d'estalviar paper, el text
   té un format pensat per ser imprès en paper DIN-A4,\index{DIN-A4} a dues cares.

    El contingut del llibre és molt divers, i va des de temes força teòrics fins a
    d'altres bastant més pràctics. He procurat donar exemples de tots els conceptes
    que s'hi expliquen, excepció feta dels molt elementals, perquè crec que és important
     no quedar-se només amb la teoria de  la resolució d'un problema determinat, sinó que
     és molt útil veure exemples resolts pas a pas.

    Encara que he fet tots els esforços possibles per eliminar qualsevol
    mena  d'error en aquest text, és gairebé inevitable que n'hagi quedat algun,
    per tant, si algú troba algun error farà bé d'avisar-me!


   Únicament em resta dir que espero que els que llegeixin aquest llibre el trobin
   útil i interessant.

\vspace*{1cm}
\hfill
\begin{minipage}[b]{35mm}
    \includegraphics{Pre-Prefaci-JMB.pdf}
\end{minipage}

{\large

\hfill \textbf{\textsl{Josep Mollera Barriga}}

\hfill \textbf{\textsl{6 de gener de 2017}}

\hfill {\small \href{mailto:josep.mollerab@outlook.com}{\faEnvelope\hspace{2mm}josep.mollerab@outlook.com} }

}
